\chapter{Conclusion and Future Work}
\label{ch:conclusion}

This thesis investigated the fundamental question of how a sequential
decision-making agent should allocate resources in a market-like environment
when the true return distribution is unknown, non-stationary, and potentially
adversarial.  Beginning from the measure-theoretic foundations of
filtered probability spaces and the Kelly Criterion, we showed that classical
Bayesian agents fail systematically when the environment changes, and we
proposed and validated two complementary remedies.

% ---------------------------------------------------------------
\section{Synthesis of Findings}
% ---------------------------------------------------------------

We demonstrated that structured exploration (Thompson Sampling) outperforms
pure exploitation (Naive Bayes Kelly) in stationary learning contexts.
However, Experiments~\ref{sec:exp4} and \ref{sec:exp5} established that
standard Bayesian updating strategies are severely vulnerable to
non-stationary environments.  The Sticky Prior Paradox
(Proposition~\ref{prop:sticky_prior}) provides a precise quantification of
this vulnerability: an agent with $T_0$ observations of past data requires
$O(T_0)$ additional observations after a regime switch before its posterior
belief shifts meaningfully, causing it to over-bet into a hostile environment
for an extended period.

To resolve these failures, we implemented and validated a two-pronged
proposed framework:

\begin{enumerate}
    \item \textbf{Volatility-Augmented HMM (Vol-HMM):} By augmenting the
    standard univariate return observation with rolling volatility and
    combining the Bayesian forward filter with a CUSUM change-point trigger,
    we demonstrated that regime detection lag drops from 14--15 steps (under
    standard Beta-posterior agents) to approximately 1.9 steps.  The
    Vol-HMM achieves a median terminal wealth of $42.3\times$ with only 3\%
    ruin risk in the harshest adversarial shock environment, compared to
    $0.78\times$ and 16\% ruin for Thompson Sampling.

    \item \textbf{Risk-Constrained CPPI:} We proved
    (Proposition~\ref{prop:cppi_bound}) that a CPPI dynamic leverage
    constraint provably bounds the maximum drawdown at a user-specified
    floor $D_{\max}$.  The CPPI achieves 0\% ruin probability across all
    synthetic environments and the highest Sortino ratio among all agents in
    the 2008 historical crisis, at the cost of lower median wealth in benign
    regimes---the quantified ``cost of survival.''
\end{enumerate}

Finally, by integrating realistic transaction costs and single-path empirical
S\&P~500 data for four distinct market epochs, we confirmed that both proposed
methods maintain their theoretical advantages under real-world conditions.
The fee audit (Section~\ref{sec:fee_audit}) further shows that the Vol-HMM's
lower turnover frequency means its net-of-fees performance advantage over
high-turnover adversarial agents is \emph{larger} than the gross advantage,
not smaller.

% ---------------------------------------------------------------
\section{Limitations}
% ---------------------------------------------------------------

Several limitations of the current framework deserve explicit acknowledgment:

\begin{enumerate}
    \item \textbf{Specified emission parameters.}  The Vol-HMM emission
    distributions are specified a priori rather than estimated via
    Baum--Welch EM (Remark~\ref{rem:hmm_specified}).  While this choice
    is justified by robustness concerns, it means the agent's performance
    depends on the accuracy of the domain-knowledge priors.  Misspecified
    priors could cause the filter to track the wrong regime.

    \item \textbf{Discrete-time CPPI floor breach.}  The analytical guarantee
    of Proposition~\ref{prop:cppi_bound} requires $m|r_{\min}|\leq 1$.  In
    discrete time, a single large overnight return (e.g., the $-20\%$ day of
    March 16, 2020) can breach the floor.  Our empirical experiments confirm
    that breaches are rare but not impossible.

    \item \textbf{Single-asset framework.}  All experiments model a single
    risky asset.  Real portfolios exhibit cross-asset correlation and
    regime coupling that the current framework does not capture.

    \item \textbf{Fractional EXP3 variant.}  Our EXP3 implementation is a
    continuous-fraction variant whose adversarial regret bound
    (Theorem~\ref{thm:exp3_regret}) does not apply exactly
    (Remark~\ref{rem:exp3_variant}).  The bound motivates the agent's design
    rather than providing a strict guarantee.
\end{enumerate}

% ---------------------------------------------------------------
\section{Future Directions}
% ---------------------------------------------------------------

Several high-impact research pathways remain open:

\begin{enumerate}
    \item \textbf{Multi-Asset Regime Coupling.}
    Extending the Vol-HMM to detect correlated regime shifts across $N$
    assets using non-Gaussian copula structures would be a natural next step.
    A joint HMM over a multivariate return vector would allow the CPPI floor
    to be defined in terms of portfolio drawdown rather than single-asset
    drawdown.

    \item \textbf{Learned Emission Parameters via Online EM.}
    Replacing the fixed emission specification with an online Baum--Welch
    update (using a sliding window to prevent over-accumulation of past data)
    would make the Vol-HMM fully adaptive, removing the dependence on domain-
    knowledge priors.

    \item \textbf{Deep Reinforcement Learning for Hyperparameter Tuning.}
    The CPPI multiplier $m$, the CUSUM threshold $h$, and the HMM persistence
    $A_{ii}$ are currently fixed by grid search.  A DRL agent trained to
    dynamically adjust these hyperparameters in response to real-time market
    entropy could further reduce detection lag and improve the growth-safety
    trade-off.

    \item \textbf{Microstructure Friction Modeling.}
    The current transaction cost model is proportional to turnover.
    Incorporating bid-ask spread decay, order book impact, and market impact
    functions would make the friction model more realistic for high-frequency
    applications and would likely further favor low-turnover agents such as
    the Vol-HMM.
\end{enumerate}

These directions represent the next iteration of protecting adaptive systems
against the inherent non-stationarity of human-centric financial markets.
